\section{Engineering Use and Scope}
The method prepares a rule for a process-aware information system. Its response card records the observation contract, action and trigger, model version, calibration range, stress scenarios, assumed effects, costs and comparison results. The process owner approves feasible actions and service requirements; an operations controller records activation, resource availability and observed consequences. Changes in workflow, staffing or case mix, persistent forecast errors, or intervention effects outside the approved range require recalibration. These integration requirements complement the evaluated computational procedure.

Policy admission must reflect organizational objectives. A lower integrated WIP cost can coexist with more sustained-overload events, as in P0 combined feedback. If a service-level constraint is binding, the engineer must exclude policies that violate it before optimizing cost. Multi-objective prescription~\cite{multikpi} addresses such trade-offs. A useful decision record therefore includes event incidence and, for shared resources, outcomes for each competing process, alongside total cost.

The evaluation covers designed synthetic structures, specified stress distributions and supplied action effects. Equal process weights do not estimate the prevalence of these mechanisms in organizations. Resource mappings and capacities are engineering inputs in the new discovery evaluation; they are not recovered automatically. Structural validity, adequate behavioral coverage and credible intervention effects are separate conditions for relying on a simulated response. Real-log prediction studies provide evidence about observable workload signal, rather than real intervention benefit.

Extrapolation remains a substantive assumption. P0 routing calibration reaches at most 19 unfinished cases, below trigger 20 and severe-event threshold 100. Logs confined to that range cannot distinguish mechanisms that agree there but diverge under stress. More complex control flow does not remove this identification limit. When a decision depends on poorly observed loads or uncertain action effects, additional measurements or a controlled pilot are needed before operational use.

\section{Conclusion}
The serial reference and structural extension show that preparing a congestion response is sensitive to both behavioral and control-flow modeling choices. Retaining load dependencies improves the serial reference and the average routing/combined comparison with supplied structure. The broader evaluation also identifies processes where it worsens decisions and discovery errors that erase the aggregate advantage. Its value depends on the process and the assumed action effects. An executable process model should therefore be assessed through prospective response outcomes, alongside prediction error and structural conformance.

\begin{credits}
\subsubsection{Artifact availability and reproducibility.}
The supplementary research archive contains protocols, code, input checksums, fitted numerical parameters, trajectories, comparator outcomes and scripts for regenerating analyses and tables. It distinguishes the primary experiment from subsequent diagnostics and includes the full supporting studies. Public event logs are identified by version-1 DOIs.
\subsubsection{Use of generative AI.}
Generative AI assisted simulator, analysis and figure-generation code, experiment orchestration, literature checking, drafting, translation and artifact-based self-review. Synthetic data were generated by the documented simulators; numerical claims derive from executed experiments and saved outputs checked by analysis scripts. Generative AI is not part of the evaluated runtime models. The authors are responsible for the manuscript and research artifacts.
\subsubsection{\discintname}
The authors have no competing interests to declare that are relevant to the content of this article.
\end{credits}
